% Part: many-valued-logic
% Chapter: sequent-calculus
% Section: rules-and-proofs

\documentclass[../../../include/open-logic-section]{subfiles}

\begin{document}

\olfileid{mvl}{seq}{rul}

\olsection{القواعد وال\usetoken{P}{derivation}}

فيما يأتي، لتدل~$\Gamma, \Delta, \Pi, \Lambda$ على متتاليات منتهية
(وربما خالية) من ال!!{sentence}s.

\begin{defn}[التتابعية]
\emph{التتابعية ذات الجوانب~$n$} تعبير من الصورة
\[
\Gamma_1 \nSequent \dots \nSequent \Gamma_n
\]
حيث كل~$\Gamma_i$ متتالية منتهية (وربما خالية) من !!{sentence}s
اللغة~$\Lang L$.
\end{defn}

\begin{defn}[التتابعية الابتدائية]
\emph{التتابعية الابتدائية ذات الجوانب~$n$} تتابعية ذات جوانب~$n$ من
الصورة~$!A \nSequent \dots \nSequent !A$ لأي !!{sentence}~$!A$ في
اللغة.

وإذا احتوت اللغة على رابط صفري الرتبة~$\star$، أي ثابت قضايا، أخذنا
أيضًا التتابعية~$\dots \nSequent \star \nSequent \dots$، حيث تظهر
$\star$ في الموضع المقابل لقيمة الصدق المرتبطة بـ
$\tf{\star} \in V$، وتكون المواضع الأخرى خالية.
\end{defn}

لكل رابط في منطق ذي~$n$ من القيم~$\Log{L}$ قاعدة منطقية لكل قيمة صدق
يمكن أن يأخذها ذلك الرابط في~$\Log{L}$. وال!!^{derivation}s في حساب
تتابعيات ذي جوانب~$n$ للمنطق~$\Log{L}$ أشجار من التتابعيات؛ تكون
التتابعيات العليا فيها تتابعيات ابتدائية، وإذا وقعت تتابعية تحت
تتابعية أخرى أو أكثر، وجب أن تتبع منها على الوجه الصحيح بقاعدة
استدلال لروابط~$\Log{L}$.

\begin{defn}[المبرهنات]
تكون !!^a{sentence}~$!A$ \emph{مبرهنة} في منطق ذي~$n$ من القيم
$\Log{L}$ إذا وُجد !!a{derivation} للتتابعية ذات الجوانب~$n$ التي
تحتوي~$!A$ في كل موضع يقابل قيمة صدق معيّنة في~$\Log{L}$. ونكتب
$\Proves[\Log{L}] !A$ إذا كانت~$!A$ مبرهنة، ونكتب
$\Proves/[\Log{L}] !A$ إذا لم تكن كذلك.
\end{defn}

\begin{defn}[!!^{derivability}]
تكون !!^a{sentence}~$!A$ \emph{!!{derivable} من} مجموعة
!!{sentence}s~$\Gamma$ في منطق ذي~$n$ من القيم~$\Log{L}$، أي
$\Gamma \Proves[\Log{L}] !A$، إذا وفقط إذا وُجدت مجموعة جزئية
منتهية~$\Gamma_0 \subseteq \Gamma$ ومتتالية~$\Gamma_0'$ من
ال!!{sentence}s في~$\Gamma_0$ بحيث يكون للتتابعية الآتية
!!a{derivation}:
\[ \Lambda_1 \nSequent \dots \nSequent \Lambda_n \]
حيث تكون~$\Lambda_i$ هي~$!A$ إذا كان الموضع~$i$ يقابل قيمة صدق
معيّنة، وتكون~$\Gamma_0'$ إذا كان يقابل قيمة صدق غير معيّنة. وإذا لم
تكن~$!A$ !!{derivable} من~$\Gamma$، كتبنا
$\Gamma \Proves/[\Log{L}] !A$.
\end{defn}

فمثلًا، لمنطق~\L ukasiewicz الثلاثي القيم حساب تتابعيات ذو ثلاثة
جوانب. وفي التتابعية الثلاثية الجوانب
$\Gamma \nSequent \Pi \nSequent \Delta$، تقابل~$\Gamma$ القيمة
$\False$، وتقابل~$\Delta$ القيمة~$\True$، وتقابل~$\Pi$ القيمة
$\Undef$. والتتابعيات الابتدائية هي
$!A \nSequent !A \nSequent !A$. ولما كانت~$\True$ وحدها معيّنة، فإن
$\Gamma \Proves[\LogLuk[3]] !A$ إذا وفقط إذا كان للتتابعية
$\Gamma \nSequent \Gamma \nSequent !A$ !!a{derivation}. (ولو كانت
$\Undef$ معيّنة أيضًا، لاحتجنا إلى !!a{derivation} للتتابعية
$\Gamma \nSequent !A \nSequent !A$.)

\end{document}
